\section{Introduction}
{\color{sub}\footnotesize 20 hrs \gap$\cdot$\gap asked in \mk{all 29 papers}
\gap$\cdot$\gap 6 syllabus units, 28 sub-topics \gap$\cdot$\gap \mk{932 marks, 39.6\% of the archive}
$\rightarrow$ the biggest chapter by far: about 21 of every 80 marks.}

\vspace{3pt}\noindent
\colorbox{gdL}{\makebox[\dimexpr\textwidth-2\fboxsep][l]{%
  \textcolor{gd}{\bfseries\footnotesize READ THIS FIRST}}}
\par\nobreak\vspace{3pt}

Chapter 1 is pure theory, and almost every question is \mk{an 8-mark question split 4+4,
3+5 or 2+6}: two lists glued together. Learn the lists below as lists.
\begin{enumerate}[label=\arabic*., leftmargin=1.4em, itemsep=1pt]
  \item \textbf{Six lists carry the chapter}, each asked in 9 to 16 papers: importance of
    organization, principles of organization, formal vs informal organization, functions of
    management, levels + skills of management, and the three classical theories (Taylor,
    Fayol, Mayo).
  \item \textbf{Forms of ownership and structures} are compare-and-choose questions:
    ``JSC vs partnership'', ``line \& staff vs line'', ``which structure for an engineering
    project''. The comparison tables in $\S$1.4 and $\S$1.5 answer them directly.
  \item \textbf{Purchasing and marketing}: two questions set every year in some form
    (functions of purchasing; importance of marketing).
\end{enumerate}
Tier chips count \textbf{papers out of 29} (2065--2082 BS):
\tS{n} asked in 8+ papers \gap \tF{n} 4--7 papers \gap \tP{n} 1--3 papers inside an
8-mark question \gap \tL{n} 1--3 papers, short note only. Bold year = Regular exam.

\hr

%=======================================================================
\T{1.1 Organization}

\Q{Define organization. Why is it necessary? Explain its importance and role in society}{\m{1}\m{2}\m{3}\m{4}\m{5}\m{8}}
{\color{sub}\footnotesize \tS{16} \yr{\textbf{82 Bh}, \textbf{80 Bh}, \textbf{78 Bh}, \textbf{76 Ch}, 76 Ash, \textbf{75 Ch}, 75 Ash, 74 Ash, \textbf{73 Ch}, 73 Shr, \textbf{72 Ch}, 72 Ka, \textbf{71 Ch}, \textbf{70 Ch}, \textbf{69 Ch}, \textbf{65 Shr}} \gap$\cdot$\gap the opener to half the chapter's questions}

\lead{Definition}
\begin{itemize}
  \item \textbf{Mooney and Reiley}: ``Organization is the form of every human association for
    the attainment of a common purpose.''
  \item Working definition: a \textbf{collection of people working together in a coordinated
    and structured way to achieve one or more goals}.
  \item As a \textbf{process}: identifying + grouping work $\rightarrow$ defining +
    delegating responsibility and authority $\rightarrow$ establishing relationships so people
    work effectively together.
  \item Three views of the same thing (AJ Sir): \textbf{goal-oriented system} (people work
    for specific goals), \textbf{social system} (people work in groups),
    \textbf{technological system} (people apply knowledge and techniques).
\end{itemize}

\sbs{%
\lead{Necessity (why we need organization)}
\begin{itemize}
  \item \textbf{Complexity of industry}: large plants, complex methods need structure.
  \item \textbf{Growing competition}: required quantity + quality, at required time, at
    minimum cost.
  \item \textbf{Optimum use of resources}: men, machines, materials, money.
  \item \textbf{Fixes authority and responsibility} $\rightarrow$ no confusion, no conflict.
  \item \textbf{Coordination} of efforts of many people toward one goal.
  \item \textbf{Reduces labour problems}; \textbf{facilitates administration}.
  \item \textbf{Stimulates creativity} and specialization.
  \item People achieve together what no individual can alone; diverse skills, more
    resources.
\end{itemize}
\textit{Andrew Carnegie}: take away the money and the works but leave the organization, and
he would rebuild within a few years. Success depends on organization.}{%
\lead{Importance / role in society}
\begin{itemize}
  \item \textbf{Produces goods and services} society needs (food, power, transport).
  \item \textbf{Employment} and income $\rightarrow$ livelihood for millions.
  \item \textbf{Development}: hospitals, schools, universities, banks, telecom, roads.
  \item \textbf{Raises living standard}: better products at lower cost.
  \item \textbf{Innovation and technology}: research, new products.
  \item \textbf{Revenue to government} (taxes) $\rightarrow$ public services.
  \item \textbf{Social welfare}: CSR, cooperatives, NGOs serve the weak.
  \item \textbf{Civilization itself} is organized effort: without organization there is
    only self-sufficiency, no division of labour, no trade.
  \item \textbf{Can we sustain without organization?} No: every need beyond bare survival
    (education, health, security) is delivered by some organization.
\end{itemize}}

\creamq{Role of organization in an employee's professional growth}{\m{4}}
{\color{sub}\footnotesize \tP{2} \yr{\textbf{81 Bh}, \textbf{75 Ch}} \gap$\cdot$\gap always paired with ``do we need informal organization?'' (below)}
\begin{itemize}
  \item \textbf{Training and development}: on-the-job training, workshops, higher study.
  \item \textbf{Career path}: promotion through levels, job rotation, wider responsibility.
  \item \textbf{Performance appraisal / merit rating} $\rightarrow$ feedback, recognition.
  \item \textbf{Mentoring}: seniors and supervisors transfer knowledge (knowledge-oriented model).
  \item \textbf{Exposure}: teamwork, projects, clients, new technology build skills.
  \item \textbf{Security + remuneration}: stable income lets an employee invest in learning.
\end{itemize}

\Q{Organization as a system: open system, social system, system approach applied to organization}{\m{2}\m{4}}
{\color{sub}\footnotesize \tP{3} \yr{\textbf{81 Bh}, 79 Ba, \textbf{70 Ch}}}

\sbsr{0.50}{%
\begin{itemize}
  \item \textbf{System} = set of interdependent parts working together toward a goal.
  \item Organization as a system:
  \begin{itemize}
    \item \textbf{Inputs}: 5 M's, men, materials, machines, money, methods (+ information).
    \item \textbf{Transformation process}: sub-systems (production, finance, marketing,
      personnel) convert inputs.
    \item \textbf{Outputs}: products, services, profit, employee satisfaction.
    \item \textbf{Feedback}: customer response, results vs plan $\rightarrow$ corrections.
    \item \textbf{Environment}: external (social, political, economic, cultural,
      technology) and internal (board, employees, culture).
  \end{itemize}
  \item \textbf{Open system}: continuously exchanges inputs and outputs with its
    environment and must adapt to changes in technology, market, customer needs, rules.
    A closed system ignores the environment and decays.
  \item \textbf{Social system}: people interact in groups; their attitudes, perceptions,
    norms shape how it works (Mayo, Barnard).
  \item \textbf{System approach applied}: every change (new technology, process) ripples
    to people, jobs, skills $\rightarrow$ plan for change, assess change forces,
    implement (Kaila). Also called a change-management plan.
\end{itemize}}{%
\figT{c1_system.png}
{\color{sub}\scriptsize Organization as an open system: inputs, processing sub-systems,
outputs, feedback, environment. Adhikari notes, p18.}}

\Q{Principles of organization}{\m{4}\m{5}\m{6}\m{8}}
{\color{sub}\footnotesize \tS{9} \yr{\textbf{82 Bh}, \textbf{80 Bh}, 76 Ash, \textbf{74 Ch}, \textbf{72 Ch}, 72 Ka, \textbf{71 Ch}, 71 Shr, 70 Asa} \gap$\cdot$\gap ``with relevant examples'' in 82 Bh}

\sbs{%
\begin{enumerate}
  \item \textbf{Unity of objectives}: clear goals for firm, each department, each post.
  \item \textbf{Division of labour / specialization}: work split into jobs matched to skill.
  \item \textbf{Departmentalization}: related activities grouped into units.
  \item \textbf{Scalar chain / chain of command}: clear line of authority top to bottom.
  \item \textbf{Unity of command}: one subordinate, one boss.
  \item \textbf{Authority = responsibility}: authority given must match duty assigned.
  \item \textbf{Span of control}: limited no.\ of subordinates per supervisor (5--6 at top,
    up to 20 at shop floor).
  \item \textbf{Delegation}: authority pushed down to where decisions are made.
\end{enumerate}}{%
\begin{enumerate}[start=9]
  \item \textbf{Centralization vs decentralization} balanced to firm size.
  \item \textbf{Coordination}: departments' work synchronized.
  \item \textbf{Simplicity}: minimum levels $\rightarrow$ easy communication.
  \item \textbf{Flexibility}: structure adapts to growth and change.
  \item \textbf{Communication}: free two-way flow, with feedback.
  \item \textbf{Continuity}: structure survives change of people.
  \item \textbf{Efficiency}: goals at minimum cost and effort.
\end{enumerate}
\textbf{Examples} (for 82 Bh): a hospital has one medical director per ward (unity of
command); a factory foreman supervises 15 workers but a CEO only 5 managers (span); a
bank sets branch targets from a national target (unity of objectives).}

\creamq{Importance of setting goals for an organization}{\m{3}}
{\color{sub}\footnotesize \tP{1} \yr{\textbf{82 Bh}}}
\begin{itemize}
  \item Gives \textbf{direction}: everyone knows what to achieve.
  \item \textbf{Basis for planning} and resource allocation.
  \item \textbf{Standard for control}: actual results measured against goals.
  \item \textbf{Motivates}: clear, achievable targets (see MBO, $\S$1.3).
  \item \textbf{Coordinates} departments toward one end; resolves conflicts.
\end{itemize}

\Q{Formal and informal organization: definitions, differences, need for informal organization}{\m{2}\m{3}\m{4}\m{5}}
{\color{sub}\footnotesize \tS{11} \yr{82 Ba, \textbf{81 Bh}, \textbf{79 Bh}, 79 Ba, \textbf{76 Ch}, \textbf{75 Ch}, \textbf{74 Ch}, 73 Shr, 71 Shr, 70 Asa, \textbf{65 Shr}} \gap$\cdot$\gap short note in 65 Shr}

\begin{itemize}
  \item \textbf{Formal organization}: structure \textbf{deliberately created by management}
    for enterprise objectives; official network of authority, responsibility and
    communication. \textbf{Chester Barnard}: ``a system of consciously coordinated activities of
    two or more persons towards a common objective''; exists when people (a) can communicate,
    (b) are willing to act, (c) share a purpose.
  \item \textbf{Informal organization}: network of personal and social relations that
    \textbf{emerges spontaneously} from friendship, language, culture, common interests.
    Barnard: ``joint personal activity without conscious common purpose, though contributing
    to joint results''. Has its own informal leader, norms and grapevine.
\end{itemize}

\vspace{2pt}
\begin{tabularx}{\textwidth}{@{}L{2.4cm}XX@{}}
\toprule
\hrow \thead{Basis} & \thead{Formal} & \thead{Informal} \\
\midrule
Origin & Deliberately created by management & Spontaneous, social and psychological forces \\
Purpose & Organization's legitimate goals & Members' social satisfaction \\
Size & Can be very large & Small groups \\
Stability & Stable, long-lived & Unstable \\
Authority & Hierarchy, flows top down & All equal; power flows from the group \\
Behaviour & Rules and regulations & Group norms, beliefs, values \\
Communication & Official chain of command & Grapevine, informal channels \\
Leadership & Vested in the manager & Chosen by the group, may out-power the manager \\
Abolition & Management can abolish it & Management cannot abolish it \\
Chart & Shown on organization chart & Never shown \\
\bottomrule
\end{tabularx}

\sbs{%
\lead{Need for informal organization (``do we need it indeed?'')}
\begin{itemize}
  \item Yes: it \textbf{fills gaps} the formal structure leaves.
  \item \textbf{Social satisfaction}, sense of belonging $\rightarrow$ morale.
  \item \textbf{Fast communication} (grapevine) across departments.
  \item \textbf{Cooperation and teamwork} beyond job descriptions.
  \item \textbf{Safety valve} for frustration; eases work pressure.
  \item \textbf{Feedback to management} on real employee mood.
  \item Risks: rumours, resistance to change, group conflict. Management should
    \textbf{understand and use it}, not suppress it (Hawthorne, $\S$1.3).
\end{itemize}}{%
\lead{Informal organization within the same family? (79 Ba)}
\begin{itemize}
  \item \textbf{Yes}. A family is itself an informal organization (no written hierarchy),
    yet sub-groups form inside it: siblings, cousins of the same age, members sharing a
    hobby or a family business.
  \item In a \textbf{family business} (single owner or partnership) the formal roles
    (owner, manager, accountant) sit alongside informal family loyalties, which often
    decide more than the formal roles do.
\end{itemize}}

\creamq{Historical development of organization}{\m{3}}
{\color{sub}\footnotesize \tP{1} \yr{\textbf{79 Bh}}}
\sbs{%
\textbf{1. Pre-industrial era}
\begin{itemize}
  \item Self-sufficiency (village economy).
  \item Commercial stage with marketable surplus: barter $\rightarrow$ money system.
  \item Means of transport (ships, railways) $\rightarrow$ wider trade.
\end{itemize}}{%
\textbf{2. Modern industrial era}
\begin{itemize}
  \item Industrial revolution (steam engine, machine tools) $\rightarrow$ factories.
  \item Corporate form: joint stock company with limited liability.
  \item Managerial revolution: scientific management (Taylor, Fayol).
  \item Mechanization, automation; economic liberalization.
  \item Globalization: multinationals; GATT (1947), ISO (1947), WTO (1995), SAFTA.
\end{itemize}}

\creamq{Characteristics / elements of an organization}{\m{4}}
{\color{sub}\footnotesize \tP{2} \yr{\textbf{73 Ch}, \textbf{68 Ch}}}
\begin{itemize}
  \item \textbf{Group of people} (two or more) $+$ \textbf{common goal}.
  \item \textbf{Division of work} and specialization.
  \item \textbf{Authority-responsibility structure} (hierarchy) and \textbf{coordination}.
  \item \textbf{Communication} network; \textbf{rules and procedures}.
  \item \textbf{Resources}: men, money, materials, machines, methods.
  \item \textbf{Environment} it serves and adapts to; \textbf{continuity} beyond members.
  \item For business operations add: profit motive, risk bearing, legal identity, customers.
\end{itemize}

\creamq{Organizational behaviour: a multidisciplinary field; how management affects it}{\m{4}}
{\color{sub}\footnotesize \tF{4} \yr{\textbf{68 Ch}, \textbf{68 Ba}, \textbf{67 Asa}, \textbf{66 Bh}} \gap$\cdot$\gap 2065--2068 papers only \gap \added}
\begin{itemize}
  \item \textbf{OB} = study of how individuals and groups behave in organizations, to improve
    performance and satisfaction.
  \item \textbf{Multidisciplinary}, it borrows from:
  \begin{itemize}
    \item \textbf{Psychology}: individual learning, motivation, perception, personality.
    \item \textbf{Sociology}: groups, roles, status, culture.
    \item \textbf{Social psychology}: attitude change, group decisions, communication.
    \item \textbf{Anthropology}: values, cultures across societies.
    \item \textbf{Political science}: power, conflict, politics in organizations.
  \end{itemize}
  \item \textbf{Management affects OB} through structure, leadership style, reward system,
    job design, communication and culture it creates: autocratic control breeds
    compliance; participation breeds commitment.
\end{itemize}

\penalty0\vspace{2pt}

%=======================================================================
\T{1.2 Management}

\Q{Define management. Explain the functions of management}{\m{2}\m{3}\m{4}\m{5}\m{8}}
{\color{sub}\footnotesize \tS{15} \yr{\textbf{82 Bh}, 81 Ba, 80 Ba, 79 Ba, \textbf{78 Bh}, \textbf{76 Ch}, 76 Ash, \textbf{75 Ch}, 75 Ash, 73 Shr, 72 Ka, 70 Asa, \textbf{69 Ch}, \textbf{68 Ch}, \textbf{66 Bh}} \gap$\cdot$\gap the most asked single list in the subject}

\begin{itemize}
  \item \textbf{Simple}: getting things done through other people.
  \item \textbf{Full}: ``a social process involving coordination of human and material resources
    through planning, organizing, staffing, leading and controlling, to accomplish stated
    objectives'' \textbf{effectively} (right goals) and \textbf{efficiently} (least resources).
  \item The \textbf{brain} of an organization; as old as civilization (Sun Tzu,
    Chanakya's Arthashastra).
\end{itemize}

\sbs{%
\lead{Functions (POSDCoRB, extended)}
\begin{enumerate}
  \item \textbf{Planning}: ``thinking before doing''; deciding in advance what, how, when,
    who (Koontz \& O'Donnell). Steps: objectives $\rightarrow$ premises $\rightarrow$
    alternatives $\rightarrow$ evaluate $\rightarrow$ select $\rightarrow$ derivative plans.
    Components: goals, objectives, policies, rules, procedures, programs, schedules, budgets.
  \item \textbf{Organizing}: identify + group activities, assign them, set authority
    relationships. Backbone of management.
  \item \textbf{Staffing}: fill posts: recruit, select, place, train, appraise, pay.
  \item \textbf{Directing}: orders, guidance, supervision; needs leadership.
\end{enumerate}}{%
\begin{enumerate}[start=5]
  \item \textbf{Motivating}: stimulate willingness to work (internal + external
    motivators); see Ch.\,3.
  \item \textbf{Coordinating}: ``orderly arrangement of group effort to provide unity of
    action'' (Mooney \& Reiley). Internal, external, vertical, horizontal.
  \item \textbf{Controlling}: set standards $\rightarrow$ measure performance
    $\rightarrow$ correct deviations. Physical, financial, budgetary control.
  \item \textbf{Communicating}: sender $\rightarrow$ message $\rightarrow$ receiver
    $\rightarrow$ feedback.
\end{enumerate}
\textbf{Fayol's five} (AJ Sir): planning, organizing, directing (commanding), coordinating,
controlling. Write these five if the question says ``basic'' functions.}

\creamq{Two principal functions and why they matter most (79 Ba)}{\m{4+2}}
\begin{itemize}
  \item Choose \textbf{planning} and \textbf{controlling}: they are the two ends of the same
    loop. Planning sets the target; controlling checks the result against it and feeds back
    into the next plan. Every other function (organizing, staffing, directing) executes the
    plan in between. No plan $\rightarrow$ nothing to control; no control $\rightarrow$ plans
    are wishes.
\end{itemize}

\Q{Levels of management: why different levels are needed, and functions at each level}{\m{2}\m{4}\m{5}\m{6}}
{\color{sub}\footnotesize \tS{9} \yr{\textbf{82 Bh}, 82 Ba, \textbf{80 Bh}, \textbf{78 Bh}, 75 Ash, \textbf{74 Ch}, \textbf{73 Ch}, \textbf{72 Ch}, \textbf{69 Ch}}}

\sbsr{0.58}{%
\begin{itemize}
  \item \textbf{Top level} (board of directors, CEO, MD, chairman):
  \begin{itemize}
    \item Sets goals, objectives, policies, budget.
    \item Main function: \textbf{decision making, planning}; cross-department responsibility.
    \item Answerable for success or failure of the whole firm.
  \end{itemize}
  \item \textbf{Middle level} (department heads, section officers):
  \begin{itemize}
    \item Implements top's policies; \textbf{organizes} resources.
    \item Link between top and lower; plans for own department.
    \item Selects, trains, motivates lower managers.
  \end{itemize}
  \item \textbf{Lower / operative level} (foreman, supervisor, inspector):
  \begin{itemize}
    \item Day-to-day supervision of workers; \textbf{directing} and on-the-spot control.
    \item Passes workers' problems upward.
  \end{itemize}
\end{itemize}}{%
\figT{c1_levels.png}
{\color{sub}\scriptsize Authority rises and responsibility widens toward the top. AJ Sir,
Ch.\,1 deck, slide 15.}}

\sbsr{0.58}{%
\lead{Why levels are needed}
\begin{itemize}
  \item \textbf{Span of control}: one person cannot direct hundreds.
  \item \textbf{Division of work}: strategic vs tactical vs operational decisions.
  \item \textbf{Clear authority and responsibility}; accountability at each step.
  \item \textbf{Communication}: orders down, reports up.
  \item \textbf{Specialization and career growth}: a ladder for managers.
\end{itemize}}{%
\figT{c1_time_functions.png}
{\color{sub}\scriptsize Share of time on each function by level: top managers plan and
organize most; first-line managers lead (direct) most. Adhikari notes, p57.}}

\Q{Managerial skills and the degree needed at each level; qualities of a good manager}{\m{1}\m{3}\m{4}\m{5}}
{\color{sub}\footnotesize \tS{10} \yr{82 Ba, \textbf{81 Bh}, 81 Ba, \textbf{78 Bh}, \textbf{76 Ch}, 75 Ash, \textbf{74 Ch}, \textbf{73 Ch}, \textbf{72 Ch}, 71 Shr} \gap$\cdot$\gap short note in 81 Bh}

\sbsr{0.52}{%
\lead{Katz's three skills}
\begin{itemize}
  \item \textbf{Technical skill}: job-specific methods, processes, tools, equipment of a
    specialized field. \mk{Most needed at lower level.}
  \item \textbf{Human (interpersonal) skill}: understand, communicate with, motivate and lead
    people and groups. \mk{Needed equally at all levels.}
  \item \textbf{Conceptual skill}: see the organization as a whole, analyse and diagnose
    situations, cause vs effect, long-range decisions. \mk{Most needed at top.}
\end{itemize}
\lead{Modern manager also needs}
\begin{itemize}
  \item \textbf{Decision-making}, \textbf{communication}, \textbf{IT / digital} skills.
  \item \textbf{Time management}, \textbf{diagnostic} skill, \textbf{design} skill
    (solving problems, not just spotting them).
\end{itemize}}{%
\figT{c1_skills_levels.png}
{\color{sub}\scriptsize Degree of each skill by level: conceptual shrinks and technical grows
going down; human skill stays wide throughout. Adhikari notes, p62.}
\lead{Qualities of a good manager}
\begin{itemize}
  \item Leadership, vision, decisiveness.
  \item Communication and listening.
  \item Integrity, fairness, self-confidence.
  \item Technical knowledge + ability to learn.
  \item Emotional stability, patience; initiative.
\end{itemize}}

\creamq{Roles of a manager (Mintzberg)}{\m{4}}
{\color{sub}\footnotesize \tP{1} \yr{76 Ash} \gap \added}
\begin{itemize}
  \item \textbf{Interpersonal}: figurehead, leader, liaison.
  \item \textbf{Informational}: monitor, disseminator, spokesperson.
  \item \textbf{Decisional}: entrepreneur, disturbance handler, resource allocator, negotiator.
  \item AJ Sir's ``interpersonal, informational, decisional skills'' are these three role groups.
\end{itemize}

\creamq{Importance of management}{\m{3}\m{4}}
{\color{sub}\footnotesize \tP{2} \yr{71 Shr, 70 Asa}}
\begin{itemize}
  \item \textbf{Achieves group goals}: directs everyone to common objectives.
  \item \textbf{Optimum use of resources}: efficient + effective; less waste.
  \item \textbf{Reduces costs} $\rightarrow$ competitive prices.
  \item \textbf{Establishes sound organization}: right person, right job.
  \item \textbf{Adapts to change}: new technology, markets.
  \item \textbf{Development of society}: ``no underdeveloped countries, only undermanaged
    ones''; jobs, better living standard.
\end{itemize}

\Q{Models of management: name them, explain any two or three with examples}{\m{3}\m{4}\m{5}}
{\color{sub}\footnotesize \tF{4} \yr{\textbf{81 Bh}, \textbf{80 Bh}, 76 Ash, \textbf{70 Ch}} \gap$\cdot$\gap team-effort model as a short note in 81 Bh}

\sbs{%
\begin{enumerate}
  \item \textbf{Hierarchical}: authority from superior, responsibility to superior. Eg army,
    civil service. Limitation: ignores leadership, charisma, informal communication.
  \item \textbf{Task-oriented}: manager does five tasks (plan, organize, direct, coordinate,
    control); a cycle, not a line. Eg a production plant.
  \item \textbf{Allocational}: manager obtains, allocates, shares, monitors and returns
    resources (money, manpower, materials), ``keeping the cup topped up''. Eg finance or
    project resource manager.
  \item \textbf{Transactional}: management as a stage in the value-added chain; adds value to
    raw resources. Eg supply chain firms.
\end{enumerate}}{%
\begin{enumerate}[start=5]
  \item \textbf{Team effort}: interlocking teams; manager as team leader stressing
    cooperation, leadership, coordination; team output exceeds sum of individuals. Eg
    software development (agile team), surgical team.
  \item \textbf{Knowledge-oriented}: manager builds knowledge by experience and diffuses it
    by teaching, coaching, training $\rightarrow$ firm learns to adapt. Eg consulting, R\&D.
  \item \textbf{Goal-concentrated}: management as control toward goals; information from
    environment triggers adjustments so the plan holds. Eg sales team with targets.
\end{enumerate}
A firm usually mixes several models.}

\creamq{Management: science or art?}{\m{3}}
{\color{sub}\footnotesize \tP{1} \yr{\textbf{78 Bh}}}
\sbs{%
\begin{itemize}
  \item \textbf{Science}: systematic body of knowledge; principles (Fayol, Taylor) tested by
    observation; cause-effect. Eg time study sets a standard output.
  \item Inexact science: deals with people, so results are not perfectly repeatable.
\end{itemize}}{%
\begin{itemize}
  \item \textbf{Art}: personal skill, creativity, judgement in applying knowledge; improves
    with practice. Eg two managers apply the same rule with different success.
  \item \textbf{Both}: science gives the principles, art applies them. Complementary.
\end{itemize}}

\creamq{Organization vs management (vs administration)}{\m{2}}
{\color{sub}\footnotesize \tP{3} \yr{79 Ba, 76 Ash, \textbf{71 Ch}}}
\par\vspace{2pt}
\begin{tabularx}{\textwidth}{@{}L{2.2cm}XX@{}}
\toprule
\hrow \thead{Basis} & \thead{Organization} & \thead{Management} \\
\midrule
Meaning & Structure: people + jobs + relations & Process: planning, organizing, directing, controlling \\
Nature & A tool / framework (the body) & The force that runs it (the brain) \\
Relation & One function of management & Uses organization to reach goals \\
Output & Defined roles, authority & Goals achieved efficiently \\
\bottomrule
\end{tabularx}
\textbf{Administration} lays down policy, arranges finance (top level, thinking);
\textbf{management} executes policy within those limits (doing).

\creamq{Which management theory is most practical today / best for organizations in Nepal?}{\m{4}\m{8}}
{\color{sub}\footnotesize \tP{2} \yr{80 Ba, 74 Ash} \gap$\cdot$\gap an opinion question: take a side and justify}
\begin{itemize}
  \item Argue for the \textbf{contingency / system approach}: no single best way; mix
    theories to fit the situation.
  \item \textbf{Why, for Nepal}:
  \begin{itemize}
    \item Mixed sector: family firms, cooperatives, public corporations, MNCs, IT startups
      $\rightarrow$ each needs a different style.
    \item Labour-intensive manufacturing (garments, cement) $\rightarrow$ Taylor's time
      study and incentives still raise output.
    \item Public enterprises (NEA, NOC) $\rightarrow$ Fayol's clear hierarchy, unity of
      command, discipline cure overlapping authority.
    \item Service and IT firms, remittance-driven youth workforce $\rightarrow$ behavioural
      approach: motivation, participation, recognition to retain staff.
    \item Political instability, frequent policy change $\rightarrow$ open-system thinking,
      watching the environment.
  \end{itemize}
  \item Conclude: \textbf{contingency approach}, drawing on all the classical tools.
\end{itemize}

\penalty0\vspace{2pt}

%=======================================================================
\T{1.3 Theories of Management}

\Q{Scientific management (Taylor): principles, experiments, features, applicability today}{\m{3}\m{4}\m{5}\m{8}}
{\color{sub}\footnotesize \tS{11} \yr{\textbf{82 Bh}, 81 Ba, \textbf{80 Bh}, 79 Ba, \textbf{75 Ch}, 75 Ash, \textbf{72 Ch}, \textbf{71 Ch}, \textbf{68 Ch}, \textbf{68 Ba}, \textbf{65 Shr}} \gap$\cdot$\gap short note in 68 Ch}

\begin{itemize}
  \item \textbf{F.W. Taylor} (1856--1915), US engineer at Midvale and Bethlehem Steel;
    ``Father of Scientific Management''. Book: \textit{Principles of Scientific Management}
    (1911). Term coined by Louis Brandeis (1910); Charles Babbage anticipated it (1832).
  \item \textbf{Definition}: ``the art of knowing exactly what is to be done and the best way
    of doing it''; replacing rule of thumb with scientific study of work.
  \item \textbf{Historical development}: rule-of-thumb factories, soldiering (deliberate slow
    work), no standards $\rightarrow$ Taylor's shop studies (1880s) $\rightarrow$ Gilbreths'
    motion study, Gantt chart $\rightarrow$ mass production (Ford).
\end{itemize}

\sbs{%
\lead{Principles}
\begin{enumerate}
  \item \textbf{Science, not rule of thumb}, for each element of work.
  \item \textbf{Scientific selection, training and development} of workers.
  \item \textbf{Cooperation}, not individualism, between management and labour.
  \item \textbf{Equal division of responsibility}: management plans, workers do.
  \item Harmony, not discord; maximum output, not restricted output.
\end{enumerate}
\lead{Techniques (features)}
\begin{itemize}
  \item \textbf{Time study}: standard time $\rightarrow$ fair day's work.
  \item \textbf{Motion study}: remove wasted motions (Gilbreth).
  \item \textbf{Fatigue study}: rest pauses.
  \item \textbf{Standardization} of tools, methods, conditions.
  \item \textbf{Differential piece-rate} wages: higher rate above standard.
  \item \textbf{Functional foremanship}: 8 specialist foremen.
  \item \textbf{Management by exception}: only deviations go to top.
\end{itemize}}{%
\lead{Experiments that support it}
\begin{itemize}
  \item \textbf{Pig-iron handling} (Bethlehem Steel, 1898): 75 labourers loading 42 kg pigs;
    by method study, rest periods and incentive pay, output rose \textbf{12.5 $\rightarrow$
    47.5 tons/day} (AJ Sir's figures: 12.7 $\rightarrow$ 48.8), wages \textbf{\$1.15
    $\rightarrow$ \$1.85}.
  \item \textbf{Shovelling}: shovel sized to a 21.5 lb load for each material; 400--600 men
    cut to 140, cost per ton halved.
  \item \textbf{Metal cutting}: 26 years, 30{,}000 tests on speeds, feeds $\rightarrow$
    high-speed steel.
\end{itemize}
\lead{Still applicable today? Yes, largely}
\begin{itemize}
  \item Assembly lines, fast food, BPO call centres use standard times and methods.
  \item Lean manufacturing, Six Sigma, industrial engineering are its heirs.
  \item Incentive pay, work study, ERP-based control.
  \item But: human needs ignored, monotony, worker resentment $\rightarrow$ combine with
    behavioural approach.
\end{itemize}}

\creamq{Scientific management vs management science}{\m{3}}
{\color{sub}\footnotesize \tP{1} \yr{\textbf{72 Ch}}}
\begin{itemize}
  \item \textbf{Scientific management} (Taylor, 1900s): improves \textbf{worker productivity}
    at shop level by time and motion study, standardization, incentive wages.
  \item \textbf{Management science} (quantitative approach, 1940s on): uses
    \textbf{mathematical models}, operations research, statistics, simulation, linear
    programming for \textbf{managerial decisions} at every level.
\end{itemize}

\creamq{Role of HR manager; how HR manager implements scientific management}{\m{4+4}}
{\color{sub}\footnotesize \tP{1} \yr{\textbf{75 Ch}}}
\sbs{%
\begin{itemize}
  \item \textbf{Significance}: manpower planning, recruitment, training, appraisal, wages,
    welfare, industrial relations, retention (see Ch.\,2).
\end{itemize}}{%
\begin{itemize}
  \item \textbf{Implementing Taylor}: job analysis and standards (time study); scientific
    selection by tests; training in the one best method; differential incentive wages;
    worker-management cooperation.
\end{itemize}}

\Q{Administrative management (Fayol): 14 principles, and are they applicable today?}{\m{4}\m{8}}
{\color{sub}\footnotesize \tS{9} \yr{\textbf{82 Bh}, 80 Ba, \textbf{79 Bh}, 79 Ba, 75 Ash, 72 Ka, 71 Shr, \textbf{70 Ch}, 70 Asa}}

\begin{itemize}
  \item \textbf{Henri Fayol} (1841--1925), French mining engineer; ``Father of Modern
    Management''. \textit{General and Industrial Management} (1916).
  \item Studied management \textbf{from the top down} (Taylor: from the shop floor up).
  \item Six activities of a firm: technical, commercial, financial, security, accounting,
    \textbf{managerial} (plan, organize, command, coordinate, control).
  \item Technical ability dominates at lower levels; managerial ability at higher levels.
\end{itemize}

\sbs{%
\lead{The 14 principles}
\begin{enumerate}
  \item \textbf{Division of work}: specialization.
  \item \textbf{Authority and responsibility}: go together.
  \item \textbf{Discipline}: obedience, respect for agreements; penalties.
  \item \textbf{Unity of command}: one boss.
  \item \textbf{Unity of direction}: one plan, one head per group of activities.
  \item \textbf{Subordination of individual interest} to general interest.
  \item \textbf{Remuneration}: fair to both.
\end{enumerate}}{%
\begin{enumerate}[start=8]
  \item \textbf{Centralization}: degree set by situation.
  \item \textbf{Scalar chain}: line of authority; ``gang plank'' for speed.
  \item \textbf{Order}: right place for everything and everyone.
  \item \textbf{Equity}: kindness + justice.
  \item \textbf{Stability of tenure}: low turnover.
  \item \textbf{Initiative}: encourage it.
  \item \textbf{Esprit de corps}: union is strength, team spirit.
\end{enumerate}}

\sbs{%
\lead{Applicable today? Yes}
\begin{itemize}
  \item Every organization chart uses scalar chain, unity of command, division of work.
  \item Fayol's functions (POCCC) are still the standard management process.
  \item Matrix structures bend unity of command; startups bend centralization. Principles
    are \textbf{flexible guidelines}, applied by situation (Fayol said so himself).
\end{itemize}}{%
\lead{Experiments / contributions around it}
\begin{itemize}
  \item Fayol: 30 years running Commentry-Fourchambault, turned a near-bankrupt mining firm
    profitable by applying these principles.
  \item \textbf{Gilbreths}: motion study, process charts, therbligs.
  \item \textbf{Gantt}: Gantt chart, planned vs actual progress.
  \item \textbf{Weber}: bureaucracy (rules, hierarchy, impersonality).
\end{itemize}}

\Q{Behavioural management approach (Mayo, Hawthorne): and how it compares with scientific and administrative approaches}{\m{3}\m{4}\m{5}\m{8}}
{\color{sub}\footnotesize \tS{11} \yr{82 Ba, \textbf{81 Bh}, 81 Ba, 80 Ba, 79 Ba, \textbf{78 Bh}, \textbf{76 Ch}, 75 Ash, 73 Shr, 71 Shr, \textbf{66 Bh}}}

\begin{itemize}
  \item \textbf{Elton Mayo} and associates (1920s--30s): organization is a \textbf{social
    system}, not only a techno-economic one. Productivity depends on human behaviour and
    social environment, not only money.
  \item Also: \textbf{Maslow} (need hierarchy), \textbf{McGregor} (X-Y), \textbf{Herzberg},
    Follett, Barnard. Human relations movement.
  \item \textbf{Elements}: the individual (differences, feelings, perceptions), the work
    group (informal organization), participative management.
\end{itemize}

\sbs{%
\lead{Hawthorne experiments (Western Electric, 1924--32)}
\begin{enumerate}
  \item \textbf{Illumination} (1924--27): output rose whether light went up \textit{or down}
    $\rightarrow$ physical factors were not the cause.
  \item \textbf{Relay assembly test room}: 6 women; rest, hours, pay varied; output rose
    almost every time $\rightarrow$ attention, group morale matter.
  \item \textbf{Interviewing programme}: 21{,}000 workers $\rightarrow$ feelings, informal
    groups shape performance.
  \item \textbf{Bank wiring room}: 14 men; group set its own output norm and enforced it
    $\rightarrow$ informal group controls output.
\end{enumerate}}{%
\lead{Conclusions}
\begin{itemize}
  \item Attention from managers raises output (Hawthorne effect).
  \item Workers are social beings; money is not the only motivator.
  \item Informal groups set norms; must be integrated with the formal organization.
  \item Satisfaction $\rightarrow$ productivity; participation and communication matter.
  \item Managers need \textbf{social skills}, not only technical ones.
\end{itemize}
\lead{Limitations}
\begin{itemize}
  \item Over-stresses the social side; satisfied workers are not always productive.
  \item Hawthorne method criticized (small samples).
\end{itemize}}

\vspace{2pt}
\begin{tabularx}{\textwidth}{@{}L{2.1cm}XXX@{}}
\toprule
\hrow \thead{Basis} & \thead{Scientific (Taylor)} & \thead{Administrative (Fayol)} & \thead{Behavioural (Mayo)} \\
\midrule
Focus & Worker, shop floor & Manager, whole firm & People and groups \\
View of person & Economic man, money motivates & Part of a structure & Social man, needs motivate \\
Aim & Max output per worker & Sound structure, principles & Satisfaction $\rightarrow$ productivity \\
Method & Time and motion study & Principles from experience & Experiments, observation \\
Informal group & Ignored & Ignored & Central \\
Contribution & Standards, incentives, efficiency & Management process, 14 principles & Motivation, participation, human relations \\
Limitation & Mechanical, monotony & Rigid, universal claims & Neglects structure, economics \\
\bottomrule
\end{tabularx}
\begin{itemize}
  \item \textbf{Advantages of behavioural over scientific} (78 Bh, 73 Shr): treats workers as
    humans; uses informal groups; participation raises commitment; lower turnover and
    conflict; quality improves with morale.
  \item \textbf{Over modern theory} (82 Ba): simpler to apply; directly improves morale and
    motivation, where quantitative models ignore people; still the base of HR practice.
    Admit modern theory integrates it.
\end{itemize}

\Q{Modern management theories: system approach and contingency approach}{\m{4}\m{5}\m{7}}
{\color{sub}\footnotesize \tF{4} \yr{82 Ba, \textbf{68 Ch}, \textbf{67 Asa}, \textbf{66 Bh}} \gap$\cdot$\gap short note in 66 Bh}

\sbs{%
\begin{itemize}
  \item After 1950; four streams:
  \begin{enumerate}
    \item \textbf{Quantitative}: models, OR, simulation, linear programming.
    \item \textbf{System}: organization = open system of interrelated sub-systems;
      inputs $\rightarrow$ process $\rightarrow$ outputs + feedback (Boulding, Johnson,
      Churchman, late 1960s).
    \item \textbf{Contingency}: ``it depends''.
    \item \textbf{Operational}: Koontz, O'Donnell: universal body of knowledge applied to
      each situation.
  \end{enumerate}
  \item Traits: systems view, dynamic, multidisciplinary, multivariable, adaptive,
    probabilistic.
\end{itemize}}{%
\lead{Contingency approach: importance}
\begin{itemize}
  \item Woodward, Fiedler, Lawrence and Lorsch.
  \item \textbf{No one best way}: the right structure, style, technique depends on the
    situation (size, technology, environment, people).
  \item Framework: environment $+$ management concepts $+$ the contingent relationship
    between them (if-then).
  \item \textbf{Importance}: practical, flexible; integrates classical, behavioural and
    quantitative tools; fits fast-changing environments; trains managers to diagnose.
  \item Weakness: no firm principles, hard to teach.
\end{itemize}}

\creamq{Management by objectives (MBO)}{\m{4}}
{\color{sub}\footnotesize \tL{1} \yr{\textbf{68 Ch}} \gap$\cdot$\gap short note \gap \added}
\begin{itemize}
  \item \textbf{Peter Drucker} (1954): superior and subordinate \textbf{jointly set
    objectives}, then review performance against them.
  \item Steps: set organizational goals $\rightarrow$ agree individual objectives
    $\rightarrow$ action plans $\rightarrow$ periodic review $\rightarrow$ appraisal and reward.
  \item Merits: clear goals, commitment, self-control, objective appraisal. Demerits: time
    consuming, favours short-term measurable goals.
\end{itemize}

\penalty0\vspace{2pt}

%=======================================================================
\T{1.4 Forms of Ownership}

\Q{Forms of ownership: list with brief explanation; choosing one; how an organization changes its form}{\m{4}\m{8}}
{\color{sub}\footnotesize \tF{4} \yr{82 Ba, \textbf{74 Ch}, 73 Shr, 71 Shr}}

\sbs{%
\begin{itemize}
  \item \textbf{Firm} = ownership unit combining men, materials, machines to produce and
    sell at a profit.
  \item \textbf{Five forms}:
  \begin{enumerate}
    \item \textbf{Single ownership}: one owner, unlimited liability.
    \item \textbf{Partnership}: 2+ partners by a deed, share profit and loss.
    \item \textbf{Joint stock company}: capital from shareholders, limited liability,
      separate legal entity.
    \item \textbf{Cooperative society}: voluntary, one member one vote, service motive.
    \item \textbf{Public corporation}: created by special act, owned by government.
  \end{enumerate}
  \item \textbf{Choice depends on}: size and nature of business, capital needed, liability
    exposure, control wanted, tax, managerial skill, continuity, cost of formation,
    government rules.
\end{itemize}}{%
\lead{Changing form of ownership (73 Shr)}
\begin{itemize}
  \item Growth needs capital and skills one owner lacks $\rightarrow$ owner takes partners.
  \item Partnership's unlimited liability and instability $\rightarrow$ convert to a
    \textbf{private limited company}.
  \item More capital needed $\rightarrow$ \textbf{public limited} via IPO.
  \item Reverse path: loss-making public enterprise $\rightarrow$ privatized.
\end{itemize}
\lead{Nepali examples \added}
\begin{itemize}
  \item \textbf{Chaudhary Group}: family trading shop $\rightarrow$ group of companies
    $\rightarrow$ listed CG / Nabil holdings.
  \item \textbf{Nepal Telecommunications Corporation} $\rightarrow$ Nepal Doorsanchar
    Company Ltd (2004) $\rightarrow$ listed public company.
  \item \textbf{Bhrikuti Paper, Bansbari Leather}: public enterprises privatized in the 1990s.
\end{itemize}}

\Q{Single ownership (sole proprietorship): features, advantages, limitations}{\m{4}\m{8}}
{\color{sub}\footnotesize \tP{3} \yr{\textbf{80 Bh}, 71 Shr, \textbf{65 Shr}} \gap$\cdot$\gap short note in 65 Shr; 80 Bh compares it with a JSC (table below)}

\sbs{%
\lead{Features}
\begin{itemize}
  \item Oldest, simplest; one person owns, finances, manages, bears all risk.
  \item Owner and business are the same in law.
  \item \textbf{Unlimited liability}: personal property liable for business debts.
  \item Suits small shops, cottage industry, handicrafts, farming.
\end{itemize}
\lead{Advantages}
\begin{itemize}
  \item Easy and cheap to form and dissolve; few legal formalities.
  \item Full control, prompt decisions, flexibility.
  \item Direct motivation: all profit to owner.
  \item Business secrecy.
  \item Personal contact with customers and employees.
  \item Little government control.
\end{itemize}}{%
\lead{Limitations}
\begin{itemize}
  \item Limited capital.
  \item Limited managerial skill (one person cannot master production, marketing,
    finance, HR).
  \item Unlimited liability.
  \item Uncertain life: ends with owner's death or illness.
  \item Cannot compete with large firms; no economies of scale.
  \item Heavy workload, stress; isolation in decisions.
\end{itemize}}

\Q{Partnership: features, types of partners, advantages and difficulties}{\m{3}\m{4}\m{5}}
{\color{sub}\footnotesize \tP{3} \yr{\textbf{81 Bh}, 74 Ash, \textbf{67 Asa}} \gap$\cdot$\gap short note in 67 Asa}

\sbs{%
\lead{Features}
\begin{itemize}
  \item Association of 2+ persons to run a business and share profit.
  \item Formed by agreement: \textbf{partnership deed} (capital, profit ratio, duties,
    dissolution).
  \item Unlimited, joint and several liability (except limited partners).
  \item Mutual agency: each partner binds the firm.
  \item No separate legal entity; restricted transfer of interest.
\end{itemize}
\lead{Types of partners}
\begin{itemize}
  \item \textbf{General / active}: manage, unlimited liability.
  \item \textbf{Limited}: liability up to capital; no say in management.
  \item \textbf{Sleeping / silent}: capital only, no management, still liable.
  \item \textbf{Nominal}: lends name only; no capital.
  \item \textbf{Minor}: under 18, benefits only.
  \item \textbf{Partner by estoppel}: implies being a partner, liable to third parties.
\end{itemize}}{%
\lead{Advantages}
\begin{itemize}
  \item Easy to form; little regulation.
  \item More capital than sole owner.
  \item Complementary skills; division of labour.
  \item Shared risk; prompt decisions (few people).
  \item Flexibility; personal interest of partners.
\end{itemize}
\lead{Difficulties / disadvantages}
\begin{itemize}
  \item Unlimited liability of at least one partner.
  \item Conflict, misunderstanding between partners $\rightarrow$ instability.
  \item Limited life: ends on death, insolvency, retirement of a partner.
  \item Mutual agency: one partner's mistake binds all.
  \item Cannot easily transfer share; capital still limited.
  \item Unequal effort vs equal sharing breeds resentment.
\end{itemize}}

\Q{Joint stock company: meaning, characteristics, types, formation in Nepal, advantages and disadvantages}{\m{2}\m{3}\m{4}\m{5}\m{6}\m{8}}
{\color{sub}\footnotesize \tS{10} \yr{81 Ba, \textbf{80 Bh}, 80 Ba, 79 Ba, \textbf{76 Ch}, 76 Ash, 75 Ash, \textbf{70 Ch}, \textbf{69 Ch}, \textbf{67 Asa}} \gap$\cdot$\gap \mk{most asked form of ownership}}

\begin{itemize}
  \item \textbf{Meaning}: voluntary association of persons, registered under the Companies
    Act, whose capital is divided into transferable \textbf{shares}. Shareholders own it; an
    elected \textbf{board of directors} manages it (elected at the AGM).
\end{itemize}

\sbs{%
\lead{Characteristics}
\begin{itemize}
  \item Created by law (registration); \textbf{separate legal entity}.
  \item \textbf{Perpetual succession}: unaffected by death of shareholders.
  \item \textbf{Limited liability}: up to share value.
  \item \textbf{Transferable shares} (freely in public company).
  \item Common seal; can sue and be sued.
  \item Ownership separate from management.
  \item Large membership, large capital.
\end{itemize}
\lead{Types (Companies Act 2063, Nepal) \added}
\begin{itemize}
  \item \textbf{Private limited}: 1 to 101 shareholders; cannot invite the public to buy
    shares; share transfer restricted; name ends ``Pvt.\ Ltd.''.
  \item \textbf{Public limited}: min 7 promoters, no upper limit; min paid-up capital
    Rs 1 crore; shares offered to public and listed; ``Ltd.''.
  \item Adhikari's ``max 50'' for private is the old Indian figure; Nepal's is 101.
\end{itemize}}{%
\lead{Formation / registration in Nepal \added}
\begin{enumerate}
  \item \textbf{Promotion}: promoters identify the idea, feasibility, name.
  \item \textbf{Name approval} from the \textbf{Office of the Company Registrar} (OCR)
    (online, \code{ocr.gov.np}).
  \item \textbf{Documents} filed:
  \begin{itemize}
    \item Application form.
    \item \textbf{Memorandum of Association}: name, objectives, registered office,
      authorized / issued / paid-up capital.
    \item \textbf{Articles of Association}: internal rules, board, meetings, share
      transfer.
    \item Promoters' citizenship (passport for foreigners); agreement among promoters;
      consent of directors.
    \item Registration fee (scaled to authorized capital).
  \end{itemize}
  \item OCR scrutinizes $\rightarrow$ \textbf{certificate of incorporation}.
  \item PAN / VAT at Inland Revenue; local ward registration.
  \item Public company only: prospectus approved by SEBON, IPO, \textbf{certificate of
    commencement of business}, listing at NEPSE.
\end{enumerate}}

\sbs{%
\lead{Advantages}
\begin{itemize}
  \item Huge capital from many investors; fresh issue possible.
  \item Limited liability $\rightarrow$ attracts investors.
  \item Transferable shares $\rightarrow$ liquidity.
  \item Stability, continuity.
  \item Professional, efficient management.
  \item Economies of large-scale production.
  \item Public confidence (audited, published accounts); easy loans.
  \item Social benefit: jobs, investment opportunity.
\end{itemize}}{%
\lead{Disadvantages}
\begin{itemize}
  \item Difficult, costly, lengthy formation; legal formalities.
  \item Lack of personal interest (salaried managers, scattered owners).
  \item Delay in decisions (board meetings).
  \item Lack of secrecy.
  \item Excessive government regulation.
  \item Directors may exploit shareholders; speculation in shares; favouritism.
  \item Impersonal labour relations.
\end{itemize}}

\vspace{2pt}
\begin{tabularx}{\textwidth}{@{}L{2.2cm}XXX@{}}
\toprule
\hrow \thead{Basis} & \thead{Single ownership} & \thead{Partnership} & \thead{Joint stock company} \\
\midrule
Owners & One & 2 or more, by deed & Many shareholders \\
Formation & Easiest, cheapest & Easy, by agreement & Registration, costly \\
Legal entity & None & None & Separate \\
Liability & Unlimited & Unlimited (joint) & Limited \\
Capital & Small & Moderate & Large \\
Management & Owner & Partners & Board of directors, professional \\
Continuity & Ends with owner & Ends with a partner & Perpetual \\
Transfer of interest & At will & Needs all partners' consent & Free (public) \\
Secrecy & Full & Fair & Little \\
Regulation & Minimum & Low & Heavy \\
\bottomrule
\end{tabularx}
{\small \textbf{Why JSC is better than partnership} (79 Ba, 76 Ch, 76 Ash): limited liability,
larger capital, perpetual life, transferable shares, professional management, public
confidence and easier loans.}

\creamq{Private limited vs public limited company}{\m{5}}
{\color{sub}\footnotesize \tP{1} \yr{\textbf{67 Asa}}}
\par\vspace{2pt}
\begin{tabularx}{\textwidth}{@{}L{2.4cm}XX@{}}
\toprule
\hrow \thead{Basis} & \thead{Private limited} & \thead{Public limited} \\
\midrule
Members (Nepal) & 1 to 101 & Min 7 promoters, no max \\
Shares to public & Not allowed & Offered to public, listed \\
Transfer of shares & Restricted & Free \\
Min capital & No statutory minimum & Paid-up Rs 1 crore \\
Start of business & On incorporation & Needs certificate of commencement \\
Disclosure & Less & Full, audited, published \\
\bottomrule
\end{tabularx}

\Q{Co-operative societies: meaning, features, types, merits and demerits}{\m{5}\m{8}}
{\color{sub}\footnotesize \tF{4} \yr{\textbf{82 Bh}, \textbf{81 Bh}, \textbf{79 Bh}, 72 Ka}}

\sbs{%
\begin{itemize}
  \item \textbf{Meaning}: voluntary association of persons on equal footing for common
    economic interest; \textbf{service motive}, not profit. Motto: ``each for all and all for
    each''.
  \item \textbf{Features}: voluntary, open membership; democratic management (\textbf{one
    member one vote}, regardless of shares); limited dividend, surplus shared as patronage
    bonus; self-help, equality; separate legal entity after registration.
  \item \textbf{Nepal}: registered under the Cooperative Act 2074; a primary cooperative
    needs at least 30 members \added.
\end{itemize}
\lead{Merits}
\begin{itemize}
  \item Easy formation; limited liability; continuity.
  \item Democratic; protects weak members.
  \item Eliminates middlemen $\rightarrow$ cheaper goods, better prices.
  \item Government support, tax concessions; reduces inequality.
\end{itemize}}{%
\lead{Types}
\begin{enumerate}
  \item \textbf{Producers'}: small producers pool raw materials, tools, sell output
    (dairy, handicraft, cottage industry).
  \item \textbf{Consumers'}: buy wholesale, sell daily goods to members cheaply.
  \item \textbf{Marketing}: farmers, artisans sell produce jointly (Eg milk cooperatives).
  \item \textbf{Housing}: members get land, build houses at low cost.
  \item \textbf{Credit / savings and credit}: pool savings, lend to members at fair rates
    (most common in Nepal).
  \item \textbf{Multipurpose}: several services together.
  \item \textbf{Farming}: small farmers pool land for mechanized farming.
\end{enumerate}
\lead{Demerits}
\begin{itemize}
  \item Limited capital; inefficient, amateur management.
  \item Lack of motivation (fixed dividend); lack of secrecy.
  \item Political interference, favouritism; misuse of funds (Nepal's savings-cooperative
    scandals).
\end{itemize}}

\Q{Public corporation: features, advantages, disadvantages; vs joint stock company}{\m{4}}
{\color{sub}\footnotesize \tP{1} \yr{81 Ba}}

\sbs{%
\begin{itemize}
  \item \textbf{Meaning}: autonomous body created by a \textbf{special act of parliament},
    wholly owned by government, for public service. Nepal: Corporation Act 2021 BS and
    individual acts. Eg Nepal Electricity Authority, Nepal Oil Corporation, Nepal Airlines
    Corporation, Nepal Food Corporation.
  \item \textbf{Features}: statutory creation; separate legal entity; government ownership;
    financial and managerial autonomy; service motive; public accountability to parliament;
    board appointed by government.
\end{itemize}
\lead{Advantages}
\begin{itemize}
  \item Autonomy + public accountability: flexibility of business, control of government.
  \item Service motive, protects public interest; strategic sectors kept national.
  \item Large finance from government; continuity; limited liability.
\end{itemize}}{%
\lead{Disadvantages}
\begin{itemize}
  \item Political interference; appointments by politics, not merit.
  \item Inefficiency, bureaucracy, losses (no profit pressure).
  \item Monopoly $\rightarrow$ poor service to customers.
  \item Autonomy often only on paper; delay in decisions.
\end{itemize}
\vspace{2pt}
{\small
\begin{tabularx}{\linewidth}{@{}L{1.6cm}XX@{}}
\toprule
\hrow \thead{Basis} & \thead{JSC} & \thead{Public corp.} \\
\midrule
Created by & Registration, Companies Act & Special act \\
Owners & Private shareholders & Government \\
Motive & Profit & Public service \\
Management & Elected board & Government-appointed board \\
Accountable to & Shareholders & Parliament, public \\
Capital & Shares & Government funds, loans \\
\bottomrule
\end{tabularx}}}
{\small Adhikari's ``Public Corporations'' slides (pp111--114: domestic / foreign / alien,
double taxation) describe the \textbf{American business corporation}, i.e.\ a JSC. The
syllabus means the government statutory corporation above.}

\creamq{As CEO of a software company, which ownership and structure would you choose?}{\m{4+4}}
{\color{sub}\footnotesize \tP{1} \yr{\textbf{74 Ch}} \gap$\cdot$\gap answer with a figure}
\sbsr{0.46}{%
\begin{itemize}
  \item \textbf{Ownership}: \textbf{private limited company}:
  \begin{itemize}
    \item Limited liability protects founders.
    \item Can bring in investors and give employees shares (ESOP).
    \item Continuity; credibility with foreign clients.
    \item Converts to public later if it grows.
  \end{itemize}
  \item \textbf{Structure}: \textbf{matrix}: functional departments (development, QA,
    design, marketing, finance) supply staff to project teams under project managers.
  \begin{itemize}
    \item Many parallel client projects.
    \item Shares scarce specialists.
    \item Fast response, skill growth.
  \end{itemize}
\end{itemize}}{%
\figT{c1_matrix_org.png}
{\color{sub}\scriptsize Matrix organization: functional managers across the top, project
managers down the side. AJ Sir, Ch.\,1 deck, slide 67.}}

\penalty0\vspace{2pt}

%=======================================================================
\T{1.5 Organizational Structure}

\Q{Organization structure: meaning, factors, how it is designed for an enterprise}{\m{2}\m{4}}
{\color{sub}\footnotesize \tF{6} \yr{\textbf{79 Bh}, \textbf{75 Ch}, \textbf{71 Ch}, 71 Shr, 70 Asa, \textbf{68 Ch}} \gap$\cdot$\gap short note in 79 Bh, 75 Ch, 71 Shr}
\sbs{%
\begin{itemize}
  \item \textbf{Meaning}: the \textbf{framework} that divides total activities into groups,
    fixes superior-subordinate relations and authority, and shows who reports to whom.
    Shown on an \textbf{organization chart}.
  \item Covers division of labour, coordination, communication, workflow, formal power.
  \item \textbf{Depends on}: size, nature of product, technology, complexity of problems,
    environment, competence of staff.
  \item \textbf{Types}: line, functional, line \& staff, committee, (matrix, project).
\end{itemize}}{%
\lead{How it is designed (steps)}
\begin{enumerate}
  \item Set objectives.
  \item Identify and list activities.
  \item Group activities into departments (by function, product, area, customer).
  \item Decide levels and span of control.
  \item Assign authority and responsibility; delegate.
  \item Set coordination and communication links.
  \item Draw the chart; review as the firm grows.
\end{enumerate}
No single best structure: each firm evolves its own.}

\Q{Line organization: features, advantages and disadvantages}{\m{4}}
{\color{sub}\footnotesize \tF{4} \yr{\textbf{76 Ch}, 76 Ash, \textbf{71 Ch}, 70 Asa}}
\sbsr{0.60}{%
\begin{itemize}
  \item Oldest, simplest; also \textbf{military} or \textbf{scalar} organization.
  \item Authority flows \textbf{vertically} top to bottom; responsibility bottom to top.
  \item Each head has full control of his department (finance, production, sales); each
    department self-contained.
  \item Suits small firms, routine work, few employees.
\end{itemize}
\vspace{2pt}
{\small
\begin{tabularx}{\linewidth}{@{}XX@{}}
\toprule
\hrow \thead{Advantages} & \thead{Disadvantages} \\
\midrule
Simple, easily understood & Overload on top executives \\
Clear authority and responsibility & No specialization \\
Strong discipline & Dictatorial, one-man show \\
Unity of command & Depends on a few key persons \\
Quick decisions & Duplication of work across departments \\
Rapid communication, easy coordination & Unsuitable for large firms \\
Economical & Scope for favouritism; poor communication upward \\
\bottomrule
\end{tabularx}}}{%
\figT{c1_line_org.png}
{\color{sub}\scriptsize Line organization: a single vertical chain from managing director to
foremen. Adhikari notes, p120.}}

\Q{Functional organization: features, advantages and disadvantages (with sketch)}{\m{3}\m{4}}
{\color{sub}\footnotesize \tF{4} \yr{\textbf{82 Bh}, \textbf{68 Ch}, \textbf{68 Ba}, \textbf{67 Asa}} \gap$\cdot$\gap short note in 67 Asa}
\sbsr{0.46}{%
\begin{itemize}
  \item Proposed by \textbf{F.W. Taylor}: all-round foremen were hard to find.
  \item Work divided by \textbf{function}; each function under a \textbf{specialist}
    who commands workers \textbf{only in his field}.
  \item Taylor's \textbf{functional foremanship}: 8 foremen, 4 in the planning room (route
    clerk, instruction card clerk, time and cost clerk, disciplinarian) + 4 on the shop
    floor (gang boss, speed boss, repair boss, inspector).
  \item Worker receives orders from several bosses $\rightarrow$ breaks unity of command.
\end{itemize}}{%
\figT{c1_functional_org.png}
{\color{sub}\scriptsize Functional structure: marketing, engineering, finance, production,
each under its expert. Adhikari notes, p124.}}
{\small
\begin{tabularx}{\textwidth}{@{}XX@{}}
\toprule
\hrow \thead{Advantages} & \thead{Disadvantages} \\
\midrule
Specialization, expert guidance & Violates unity of command $\rightarrow$ indiscipline \\
Efficiency, standardized operations & Overlapping authority, conflicts \\
Easy selection and training of specialists & Shifting of responsibility \\
Reduced burden on top executives & Poor coordination between functions \\
Lower cost through expertise; scope for growth & Delayed decisions needing several experts; costly \\
\bottomrule
\end{tabularx}}

\Q{Line and staff organization; its advantages over line and functional organization}{\m{4}\m{5}}
{\color{sub}\footnotesize \tF{4} \yr{82 Ba, \textbf{79 Bh}, \textbf{74 Ch}, \textbf{66 Bh}}}
\sbsr{0.50}{%
\begin{itemize}
  \item Line structure $+$ \textbf{staff specialists} (legal adviser, personnel, research,
    quality control) who \textbf{advise} line managers but have \textbf{no command
    authority}.
  \item ``Staff are thinkers, line are doers.''
  \item Suits large plants: steel mills, electricity boards, big factories.
\end{itemize}
\lead{Advantages over line organization}
\begin{itemize}
  \item Expert advice $\rightarrow$ better decisions.
  \item Relieves overloaded line executives.
  \item Suits growth and complexity.
\end{itemize}
\lead{Advantages over functional organization}
\begin{itemize}
  \item Keeps unity of command and discipline.
  \item No overlapping authority.
  \item Clear accountability for results.
\end{itemize}}{%
\figT{c1_line_staff_org.png}
{\color{sub}\scriptsize Line and staff: dashed lines are staff (advice only). Adhikari notes,
p128.}
\lead{Disadvantages}
\begin{itemize}
  \item Line-staff friction; line resents advice.
  \item Staff without authority may be ineffective.
  \item Staff not accountable for results $\rightarrow$ lethargy.
  \item Expensive; confusion over duties.
  \item Line may lose initiative, over-rely on staff.
\end{itemize}}

\Q{Committee organization: types, merits and demerits}{\m{4}}
{\color{sub}\footnotesize \tP{3} \yr{\textbf{79 Bh}, 75 Ash, \textbf{74 Ch}}}
\sbs{%
\begin{itemize}
  \item Formal group of 2+ persons appointed to deliberate on problems one person cannot
    solve; members meet, discuss, recommend or decide.
  \item Not a separate structure: added to line or line \& staff at any level.
\end{itemize}
\lead{Types}
\begin{itemize}
  \item \textbf{Ad hoc}: temporary, dissolves after its task.
  \item \textbf{Standing / permanent}: repetitive work (purchase, research committee).
  \item \textbf{Advisory (staff)}: recommend only.
  \item \textbf{Executive (line)}: has power to decide and enforce.
  \item \textbf{Formal / informal}: by charter or by convenience.
\end{itemize}}{%
\lead{Merits}
\begin{itemize}
  \item Pooled knowledge; problem seen from several views.
  \item Coordination between departments.
  \item Democratic, participative; better acceptance.
  \item New ideas; checks concentration of power.
\end{itemize}
\lead{Demerits}
\begin{itemize}
  \item Slow decisions; costly (many people's time).
  \item Compromise, not the best decision.
  \item Divided responsibility (no one accountable).
  \item Dominated by a few members.
\end{itemize}}

\Q{Matrix / project organization; which structure suits a (temporary) engineering project?}{\m{5}\m{8}}
{\color{sub}\footnotesize \tP{3} \yr{\textbf{73 Ch}, \textbf{72 Ch}, \textbf{68 Ba}} \gap$\cdot$\gap 68 Ba asks the matrix chart; 74 Ch's software-company version is in $\S$1.4}
\sbs{%
\begin{itemize}
  \item \textbf{Project organization}: a temporary team under a \textbf{project manager},
    formed for one project, dissolved at completion.
  \item \textbf{Matrix}: project structure laid over a functional one; staff report to a
    functional manager (technical) \textbf{and} a project manager (schedule, cost). Two
    bosses. (Figure in $\S$1.4.)
  \item \textbf{Merits}: efficient use of specialists across projects; clear project
    accountability; skill development, job rotation; fast response.
  \item \textbf{Demerits}: violates unity of command; power struggle; complex, confusing;
    costly coordination.
\end{itemize}}{%
\lead{Best for a temporary engineering project: project / matrix}
\begin{itemize}
  \item Project has fixed \textbf{start, end, budget, scope}: a dedicated project manager owns
    all three.
  \item Needs \textbf{many disciplines} (civil, electrical, mechanical, finance) together.
  \item Team \textbf{disbands} at completion; specialists return to departments.
  \item Single point of contact for client; faster decisions than line.
  \item \textbf{Line} would lack specialists, \textbf{functional} would scatter
    responsibility, \textbf{committees} are too slow.
  \item Large one-off project (hydropower, highway) $\rightarrow$ pure project organization;
    firm running many projects $\rightarrow$ matrix.
\end{itemize}}

\Q{Authority and responsibility; span of control; policy group vs executive group}{\m{3}\m{4}\m{5}}
{\color{sub}\footnotesize \tF{4} \yr{\textbf{68 Ch}, \textbf{68 Ba}, \textbf{67 Asa}, \textbf{65 Shr}} \gap$\cdot$\gap span of control also a short note in 68 Ch}
\sbs{%
\lead{Authority vs responsibility}
{\small
\begin{tabularx}{\linewidth}{@{}L{1.5cm}XX@{}}
\toprule
\hrow \thead{Basis} & \thead{Authority} & \thead{Responsibility} \\
\midrule
Meaning & Right to decide, command, use resources & Obligation to perform assigned work \\
Flow & Downward & Upward \\
Delegation & Can be delegated & Cannot be delegated fully \\
Source & Position (formal) & Superior-subordinate relation \\
\bottomrule
\end{tabularx}}
Principle: authority must be \textbf{equal} to responsibility.
}{%
\lead{Policy group vs executive group}
\begin{itemize}
  \item \textbf{Policy group}: board of directors / top management; frames objectives,
    policies, budgets (administration, thinking).
  \item \textbf{Executive group}: managers who carry policies out, direct day-to-day work
    (management, doing).
\end{itemize}}

\sbs{%
\lead{Span of control}
\begin{itemize}
  \item \textbf{No.\ of subordinates directly reporting} to one superior.
  \item \textbf{Narrow span}: close supervision, many levels, tall structure, slow
    communication.
  \item \textbf{Wide span}: flat structure, fast communication, more delegation, less
    control.
  \item Typical: 5--6 at top, 15--20 at shop floor.
  \item In a \textbf{line organization} span decides the no.\ of levels.
\end{itemize}}{%
\begin{itemize}
  \item \textbf{Factors}: ability of manager and subordinates, nature of work (routine =
    wide), degree of delegation, communication tools, geographic spread.
  \item \textbf{Graicunas}: with $n$ subordinates the relationships to manage are
    $n\,(2^{n-1} + n - 1)$: 6 people $\rightarrow$ 222. Why span must stay limited.
\end{itemize}}

\penalty0\vspace{2pt}

%=======================================================================
\T{1.6 Purchasing and Marketing Management}

\Q{Purchasing: meaning, objectives (5 R's), principles, functions of the purchasing department}{\m{3}\m{4}\m{5}\m{8}}
{\color{sub}\footnotesize \tS{10} \yr{\textbf{82 Bh}, \textbf{81 Bh}, \textbf{78 Bh}, \textbf{76 Ch}, 72 Ka, \textbf{71 Ch}, \textbf{70 Ch}, 70 Asa, \textbf{67 Asa}, \textbf{65 Shr}} \gap$\cdot$\gap short note ``role of purchasing and marketing'' in 82 Bh}

\sbs{%
\begin{itemize}
  \item \textbf{Purchasing}: acquiring materials, supplies, machines, tools and services
    by payment, for the firm's operation. First phase of materials management.
  \item \textbf{Alford and Beatty}: ``procuring of materials, supplies, machines, tools and
    services required for equipment, maintenance and operation of a manufacturing plant.''
  \item \textbf{Objectives: the 5 R's}: right \textbf{quality}, right \textbf{quantity},
    right \textbf{price}, right \textbf{time}, right \textbf{source} (+ right place).
  \item Continuous supply, minimum inventory, good vendor relations, low total cost.
\end{itemize}
\lead{Principles of purchasing}
\begin{itemize}
  \item Buy to specification and standards.
  \item Competitive quotations; buy on total value, not lowest price.
  \item Centralized purchasing for economy.
  \item Ethical dealing, no favouritism; prompt payment.
  \item Develop reliable, alternative sources.
  \item Keep records; review make-or-buy.
\end{itemize}
\lead{Purchasing vs procurement (71 Ch)}
\begin{itemize}
  \item \textbf{Purchasing}: the buying transaction (order, pay).
  \item \textbf{Procurement}: wider: need identification, sourcing, contracts, buying,
    logistics, receiving, vendor management; includes lease, hire, transfer.
\end{itemize}}{%
\lead{Functions of purchasing department}
\begin{enumerate}
  \item Receive and check purchase requisitions.
  \item Review specifications; standardize, simplify.
  \item Keep records of materials, substitutes, suppliers, prices.
  \item Call and analyse quotations; select vendor.
  \item Negotiate; place orders; follow up.
  \item Inspect: right quality and quantity.
  \item Prepare purchase budget.
  \item Buy at right time and economical rate.
  \item Keep uninterrupted supply to production.
  \item Liaison between vendors and production, QC, finance, stores.
  \item Ensure prompt payment $\rightarrow$ good vendor relations.
  \item Develop new sources; handle sub-contracts.
  \item Keep inventory optimum; dispose of scrap and salvage.
\end{enumerate}
\lead{Role of purchasing and marketing departments}
\begin{itemize}
  \item \textbf{Purchasing}: input side, supplies right materials at lowest total cost
    $\rightarrow$ cost and quality of product.
  \item \textbf{Marketing}: output side, finds needs, sells product, brings revenue.
  \item Together they link the firm to suppliers and customers.
\end{itemize}}

\Q{Methods of purchasing; purchasing procedure (process outline)}{\m{2}\m{4}\m{5}}
{\color{sub}\footnotesize \tF{6} \yr{\textbf{81 Bh}, 79 Ba, 76 Ash, 74 Ash, \textbf{72 Ch}, 71 Shr}}
\sbs{%
\lead{Methods}
\begin{enumerate}
  \item \textbf{By requirement} (hand-to-mouth): buy as need arises.
  \item \textbf{For a specified period}: regular items bought for a schedule.
  \item \textbf{Market purchasing}: buy ahead when prices are low and expected to rise.
  \item \textbf{Speculative}: buy excess for profit from price rise.
  \item \textbf{Rate contract}: agreed price for 2--3 years (coal, steel).
  \item \textbf{Group purchasing}: many items in one order.
  \item \textbf{Through government supplies agency} at rate-contract prices.
\end{enumerate}
\lead{By procedure (AJ Sir)}
\begin{itemize}
  \item \textbf{Direct}: urgent, from a known supplier.
  \item \textbf{Quotation}: at least 3 quotations, pick best.
  \item \textbf{Tender}: public notice in newspapers; sealed bids; evaluate (large or
    public purchases; Public Procurement Act 2063 in Nepal).
\end{itemize}}{%
\lead{Purchasing procedure (steps)}
\begin{enumerate}
  \item Recognize need: \textbf{purchase requisition} from department / store.
  \item Describe need: specification, quantity.
  \item Enquiry to suppliers.
  \item Receive quotations / tenders.
  \item \textbf{Comparative statement}; evaluate.
  \item Select and approve supplier; negotiate.
  \item \textbf{Place purchase order}; copies to stores, accounts, requesting department.
  \item \textbf{Follow up} the order.
  \item Receive goods; \textbf{inspect}; goods-received report.
  \item Check invoice; pass to accounts for \textbf{payment}.
  \item Issue materials; maintain records.
\end{enumerate}
{\color{sub}\footnotesize Draw steps 1--11 as a vertical flow of boxes for ``draw an outline
of the purchasing process'' (76 Ash).}}

\Q{Marketing: definition, concept, and its importance (in modern business and the digital era)}{\m{3}\m{4}\m{5}\m{8}}
{\color{sub}\footnotesize \tS{14} \yr{\textbf{82 Bh}, 82 Ba, \textbf{81 Bh}, 81 Ba, \textbf{79 Bh}, 79 Ba, \textbf{78 Bh}, \textbf{76 Ch}, \textbf{75 Ch}, 74 Ash, \textbf{73 Ch}, 73 Shr, \textbf{70 Ch}, \textbf{68 Ba}} \gap$\cdot$\gap \mk{asked almost every year}}

\sbs{%
\begin{itemize}
  \item \textbf{Stanton}: ``a total system of interacting business activities designed to
    plan, price, promote and distribute want-satisfying products and services to present
    and potential customers.''
  \item \textbf{Kotler}: satisfying needs and wants through exchange processes.
  \item Starts with identifying a consumer need, ends with satisfying it.
  \item \textbf{4 P's}: product, price, place, promotion.
\end{itemize}
\lead{Marketing concept}
\begin{itemize}
  \item Evolution: production $\rightarrow$ product $\rightarrow$ selling $\rightarrow$
    \textbf{marketing} $\rightarrow$ societal concept.
  \item Three pillars: \textbf{customer orientation}, \textbf{profit through satisfaction}
    (not volume), \textbf{integrated effort} of all departments.
  \item Selling concept: ``sell what we make''. Marketing concept: ``make what customers want''.
\end{itemize}
\lead{Marketing research (4 elements)}
\begin{itemize}
  \item Current sales analysis; market intelligence (competitors); market analysis
    (reaction to price, new product); product evaluation.
\end{itemize}}{%
\lead{Importance}
\begin{itemize}
  \item \textbf{To the firm}: brings revenue, the only source of profit; tells production
    what to make; builds brand, loyalty; survives competition.
  \item \textbf{To consumers}: right product, right place, right price; choice,
    information.
  \item \textbf{To society / economy}: creates demand $\rightarrow$ production $\rightarrow$
    jobs; raises living standard; distribution to remote areas; exports.
  \item Fundamental: without a customer nothing else in the firm is needed (Drucker:
    ``marketing and innovation are the two basic functions'').
\end{itemize}
\lead{Importance in the digital era (81 Ba)}
\begin{itemize}
  \item Social media, search, e-commerce: reach millions at low cost.
  \item Data analytics: precise targeting, personalization.
  \item Measurable ROI; real-time campaigns.
  \item Two-way engagement, reviews, influencers build trust.
  \item Small firms and startups compete globally.
  \item Digital payments (eSewa, Khalti) enable online selling in Nepal.
\end{itemize}}

\creamq{Functions of marketing}{\m{4}\m{5}}
{\color{sub}\footnotesize \tP{3} \yr{81 Ba, 71 Shr, \textbf{71 Ch}}}
\sbs{%
\textbf{Classical grouping} (Adhikari)
\begin{itemize}
  \item \textbf{Exchange}: buying, selling.
  \item \textbf{Physical distribution}: transport, storage / warehousing.
  \item \textbf{Facilitating}: standardization and grading, financing, risk bearing, market
    information.
\end{itemize}}{%
\textbf{Modern list} (AJ Sir)
\begin{itemize}
  \item Product / service management, pricing, promotion.
  \item Distribution (channel), selling, financing.
  \item Market information management, buying.
  \item Packaging, branding, after-sales service.
\end{itemize}}

\creamq{Advertising: definition, and why it is one of the best forms of marketing}{\m{1}\m{3}}
{\color{sub}\footnotesize \tP{3} \yr{\textbf{78 Bh}, \textbf{72 Ch}, \textbf{70 Ch}}}
\begin{itemize}
  \item \textbf{Definition}: paid, non-personal communication of information about a product,
    service or idea through mass media (newspaper, TV, radio, online, billboards) by an
    identified sponsor, to create or maintain demand. From Latin \textit{ad verto}, ``to turn
    towards''.
  \item \textbf{Why the best form}:
  \begin{itemize}
    \item Mass reach at low cost per contact; repetition builds memory.
    \item Creates awareness and new demand; launches new products.
    \item Builds brand image and goodwill; supports salesmen and dealers.
    \item Informs and educates consumers; counters competition.
    \item Digital ads are targeted and measurable.
  \end{itemize}
\end{itemize}

\creamq{Salesmanship; marketing vs purchasing; marketing software products in Nepal}{\m{4}}
{\color{sub}\footnotesize \tP{3} \yr{76 Ash, 73 Shr, \textbf{68 Ch}}}
\sbs{%
\lead{Salesmanship is an important ingredient of marketing? Agree (73 Shr)}
\begin{itemize}
  \item Personal selling persuades face to face, explains technical products, handles
    objections, \textbf{closes the sale}.
  \item Brings customer feedback to the firm; builds relationships.
  \item But it is \textbf{one ingredient}: product, price, place and research come first.
\end{itemize}
\lead{Marketing vs purchasing (68 Ch)}
\begin{itemize}
  \item Purchasing: \textbf{inputs}, firm is buyer, aims at low cost, deals with suppliers.
  \item Marketing: \textbf{outputs}, firm is seller, aims at revenue and satisfaction,
    deals with customers.
\end{itemize}}{%
\lead{Challenges for marketing software in Nepal (76 Ash)}
\begin{itemize}
  \item Small domestic market; low willingness to pay for software.
  \item Piracy; free alternatives.
  \item Preference for foreign brands; trust deficit in local products.
  \item Low digital literacy and internet access outside cities.
  \item Limits on international payments; weak IP law enforcement.
  \item Skilled marketers scarce; brain drain.
  \item Response: freemium, local language, SaaS subscriptions, export markets,
    partnerships with banks and telecoms.
\end{itemize}}

\penalty0\vspace{2pt}

%=======================================================================
\T{1.7 Topics Set Only by the 2065--2068 Papers}

{\color{sub}\footnotesize These came from the older ME708 syllabus and have not been set since
2068. Know one paragraph each. \added}

\creamq{Activities of production development}{\m{5}}
{\color{sub}\footnotesize \tP{2} \yr{\textbf{66 Bh}, \textbf{65 Shr}}}
\begin{itemize}
  \item Idea generation (market, R\&D, competitors).
  \item Screening; feasibility (technical, economic).
  \item Product design; prototype; testing.
  \item Pilot production; test marketing.
  \item Commercial launch; review.
  \item Supporting: standardization, simplification, specialization, value analysis.
\end{itemize}

\creamq{Manufacturing methods; CIM}{\m{4}\m{6}}
{\color{sub}\footnotesize \tP{1} \yr{\textbf{67 Asa}} \gap$\cdot$\gap CIM short note}
\begin{itemize}
  \item \textbf{Job / project}: one-off to order (ships, dams).
  \item \textbf{Batch}: lots of similar items (medicine, garments).
  \item \textbf{Mass / flow}: standard products on lines (cars).
  \item \textbf{Continuous process}: round-the-clock flow (cement, refinery).
  \item \textbf{CIM}: whole plant computer-controlled: CAD, CAPP, CAM, CNC, robots, FMS,
    AGVs, shared database $\rightarrow$ flexibility, quality, short lead time; high cost.
\end{itemize}

\penalty0\vspace{2pt}

%=======================================================================
\T{1.8 PYQ Mapping}

\lead{Theory questions, covered in this document}
\begin{itemize}
  \item \textbf{Define organization; need, importance, role in society; can we sustain
    without it} $\rightarrow$ $\S$1.1. 16 papers; the rider on most 8-mark questions.
  \item \textbf{Organization as open / social system; system approach} $\rightarrow$
    $\S$1.1. \yr{\textbf{81 Bh} \m{2}, 79 Ba \m{2+2}, \textbf{70 Ch} \m{4}}
  \item \textbf{Role in employee's professional growth} $\rightarrow$ $\S$1.1.
    \yr{\textbf{81 Bh} \m{4}, \textbf{75 Ch} \m{4}}
  \item \textbf{Principles of organization} $\rightarrow$ $\S$1.1. 9 papers,
    \m{4} to \m{8}.
  \item \textbf{Formal vs informal; need for informal; informal organization in a family}
    $\rightarrow$ $\S$1.1. 11 papers.
  \item \textbf{Historical development} \yr{\textbf{79 Bh} \m{3}}; \textbf{characteristics / elements}
    \yr{\textbf{73 Ch} \m{4}, \textbf{68 Ch} \m{4}}; \textbf{organizational behaviour}
    \yr{\textbf{68 Ch} \m{4}, \textbf{68 Ba} \m{4}, \textbf{67 Asa} \m{4}, \textbf{66 Bh} \m{4}} $\rightarrow$ $\S$1.1.
  \item \textbf{Functions of management} (incl.\ ``two principal functions'', 79 Ba)
    $\rightarrow$ $\S$1.2. 15 papers: the most asked list in the subject.
  \item \textbf{Levels, need for levels, functions at each level} $\rightarrow$ $\S$1.2.
    9 papers.
  \item \textbf{Managerial skills, degree per level, qualities of a manager} $\rightarrow$
    $\S$1.2. 10 papers.
  \item \textbf{Models of management} $\rightarrow$ $\S$1.2. \yr{\textbf{81 Bh} \m{4}, \textbf{80 Bh} \m{4},
    76 Ash \m{4}, \textbf{70 Ch} \m{3+5}}
  \item \textbf{Roles of a manager} \yr{76 Ash \m{4}}; \textbf{importance of management}
    \yr{71 Shr \m{4}, 70 Asa \m{3}}; \textbf{science or art} \yr{\textbf{78 Bh} \m{3}};
    \textbf{organization vs management} \yr{79 Ba \m{2}, 76 Ash \m{2}, \textbf{71 Ch} \m{2}};
    \textbf{best theory today / for Nepal} \yr{80 Ba \m{4}, 74 Ash \m{8}} $\rightarrow$
    $\S$1.2.
  \item \textbf{Scientific management} (principles, experiments, still applicable; vs
    management science; HR manager implementing it) $\rightarrow$ $\S$1.3. 11 papers.
  \item \textbf{Fayol's administrative theory, 14 principles, applicability} $\rightarrow$
    $\S$1.3. 9 papers.
  \item \textbf{Behavioural approach; its advantages over scientific / modern; comparisons of
    the three approaches} $\rightarrow$ $\S$1.3. 11 papers.
  \item \textbf{System and contingency approach} \yr{82 Ba \m{8}, \textbf{68 Ch} \m{5}, \textbf{67 Asa} \m{7},
    \textbf{66 Bh} \m{4}}; \textbf{MBO} \yr{\textbf{68 Ch} \m{4}} $\rightarrow$ $\S$1.3.
  \item \textbf{Forms of ownership; changing form; CEO of a software company}
    $\rightarrow$ $\S$1.4. \yr{82 Ba \m{4}, \textbf{74 Ch} \m{4+4}, 73 Shr \m{4+4}, 71 Shr \m{4}}
  \item \textbf{Single ownership} \yr{\textbf{80 Bh} \m{8}, 71 Shr \m{4}, \textbf{65 Shr} \m{4}};
    \textbf{partnership} \yr{\textbf{81 Bh} \m{3}, 74 Ash \m{3+5}, \textbf{67 Asa} \m{4}} $\rightarrow$
    $\S$1.4.
  \item \textbf{Joint stock company} (features, formation in Nepal, merits, vs partnership,
    vs public corporation, private vs public) $\rightarrow$ $\S$1.4. 10 papers.
  \item \textbf{Cooperatives} \yr{\textbf{82 Bh} \m{5}, \textbf{81 Bh} \m{5}, \textbf{79 Bh} \m{8}, 72 Ka \m{8}};
    \textbf{public corporation} \yr{81 Ba \m{4}} $\rightarrow$ $\S$1.4.
  \item \textbf{Organization structure} $\rightarrow$ $\S$1.5, including short notes in
    \yr{\textbf{79 Bh}, \textbf{75 Ch}, 71 Shr}.
  \item \textbf{Line} \yr{\textbf{76 Ch} \m{4}, 76 Ash \m{4}, \textbf{71 Ch} \m{4}, 70 Asa \m{4}};
    \textbf{functional} \yr{\textbf{82 Bh} \m{3}, \textbf{68 Ch} \m{3}, \textbf{68 Ba} \m{4}, \textbf{67 Asa} \m{4}};
    \textbf{line \& staff} \yr{82 Ba \m{4}, \textbf{79 Bh} \m{4}, \textbf{74 Ch} \m{4}, \textbf{66 Bh} \m{5}};
    \textbf{committee} \yr{\textbf{79 Bh} \m{4}, 75 Ash \m{4}, \textbf{74 Ch} \m{4}} $\rightarrow$ $\S$1.5.
  \item \textbf{Matrix; structure for an engineering project} \yr{\textbf{73 Ch} \m{8}, \textbf{72 Ch} \m{8},
    \textbf{68 Ba} \m{5}}; \textbf{authority, responsibility, span, policy vs executive group}
    \yr{\textbf{68 Ch}, \textbf{68 Ba}, \textbf{67 Asa}, \textbf{65 Shr}} $\rightarrow$ $\S$1.5.
  \item \textbf{Purchasing, 5 R's, principles, functions} $\rightarrow$ $\S$1.6. 10 papers.
  \item \textbf{Methods of purchasing; purchasing procedure} $\rightarrow$ $\S$1.6. 6 papers.
  \item \textbf{Marketing: definition, concept, importance} $\rightarrow$ $\S$1.6. 14 papers.
  \item \textbf{Functions of marketing; advertising; salesmanship; marketing vs purchasing;
    software marketing in Nepal} $\rightarrow$ $\S$1.6.
  \item \textbf{Production development, manufacturing methods, CIM} $\rightarrow$ $\S$1.7.
    \yr{\textbf{67 Asa}, \textbf{66 Bh}, \textbf{65 Shr}}
\end{itemize}

\lead{Short-note prompts from the Detailed PYQ's ``Extra'' section, placed here}
\begin{itemize}
  \item Scientific management, span of control, MBO \yr{\textbf{68 Ch}}; contingency management
    theory \yr{\textbf{66 Bh}}; marketing concept \yr{\textbf{68 Ba}}; partnership, functional organization,
    CIM \yr{\textbf{67 Asa}}; single ownership, informal organization \yr{\textbf{65 Shr}}; team effort model,
    managerial skills \yr{\textbf{81 Bh}}; organizational structure \yr{\textbf{79 Bh}, \textbf{75 Ch}, 71 Shr}.
  \item Industrial relations, collective bargaining, job design, incentive programs and
    manpower planning go to Ch.\,2; leadership style, contingency approach to leadership,
    satisfaction-progression vs frustration-regression and entrepreneurial characteristics
    go to Ch.\,3.
\end{itemize}

\lead{Numerical content}
\begin{itemize}
  \item \textbf{This chapter has no numerical component}: no paper in 29 sets a calculation.
    The only formula here (Graicunas' span relationships) is there to strengthen a theory
    answer.
\end{itemize}

\lead{Corrections to the source notes}
\begin{itemize}
  \item Adhikari p17: WTO was founded in 1995 (not 1955); GATT was signed in 1947.
  \item Adhikari p105: a private company in Nepal may have up to 101 shareholders (Companies
    Act 2063), not 50; p106 says ``maximum number of persons to form a company is seven'':
    seven is the \textbf{minimum} for a public company.
  \item Adhikari pp111--114: ``public corporation'' there is the US business corporation;
    the syllabus means a government statutory corporation.
  \item AJ Sir slide 50: cooperatives ``at least 10 and no more than 100 members'' is not
    Nepal law; the Cooperative Act 2074 sets a minimum of 30 for a primary cooperative.
  \item AJ Sir slide 31: Fayol was born in 1841.
  \item Syllabus file prints course code CT 710; every paper prints ME708.
\end{itemize}
